\section{Marco Teórico}

El desarrollo de soluciones automatizadas para la prevención del hurto en el comercio minorista requiere un marco conceptual integrado que transite, de manera lógica y jerárquica, desde los fundamentos de la videovigilancia computacional hasta las arquitecturas algorítmicas de aprendizaje profundo y sus métricas de evaluación operacional en dispositivos de cómputo limitado.

\subsection{Videovigilancia y Seguridad en el Entorno Retail}

\subsubsection{Videovigilancia Inteligente (\textit{Intelligent Video Surveillance})}

Los sistemas tradicionales de Circuito Cerrado de Televisión (CCTV) han operado históricamente como mecanismos pasivos de registro, cuya eficacia depende críticamente de la supervisión humana en tiempo real. Sin embargo, la fatiga visual, la saturación de monitores y la limitación de la atención selectiva reducen drásticamente la capacidad operativa de los guardias de seguridad para detectar eventos anómalos en el momento en que ocurren.

Frente a estas deficiencias, la Videovigilancia Inteligente (\textit{Intelligent Video Surveillance} -- IVS) se constituye como el paradigma orientado a transformar los flujos de video continuos en información estructurada mediante el procesamiento analítico autónomo \cite{anikwe-onyiwa_recent_2025}.

La función primordial de un sistema IVS reside en la automatización del monitoreo mediante algoritmos capaces de interpretar escenas complejas sin intervención humana constante, transitando de una seguridad puramente reactiva, donde la grabación se consulta únicamente tras la consumación de una pérdida, hacia una prevención proactiva en tiempo real.

En el sector \textit{retail}, la implementación de la IVS busca mitigar la merma desconocida (\textit{inventory shrinkage}), la cual engloba las pérdidas de inventario derivadas del hurto sin violencia. La característica distintiva de las soluciones modernas de IVS es su capacidad para operar de forma continua, filtrando el ruido visual de las actividades rutinarias de los clientes y aislando patrones de movimiento relevantes para la seguridad del establecimiento \cite{hameed_real_2025}.

\subsubsection{Reconocimiento de Actividades Humanas (\textit{Human Activity Recognition})}

Dentro del dominio de la visión por computadora, el Reconocimiento de Actividades Humanas (\textit{Human Activity Recognition} -- HAR) constituye la disciplina encargada de identificar, clasificar e interpretar las acciones ejecutadas por uno o más individuos a partir de secuencias de imágenes o flujos de video.

Mientras que la detección de objetos se restringe a responder qué elementos están presentes en la escena y dónde se ubican, el HAR aborda la dimensión dinámica, respondiendo a la interrogante de qué acción está realizando el sujeto a lo largo del tiempo \cite{ansari_optimized_2023}.

Las actividades humanas se categorizan comúnmente según su nivel de abstracción temporal en tres niveles: gestos (movimientos atómicos de extremidades), acciones (secuencias simples como caminar, inclinarse o tomar un producto) y actividades complejas (cadenas de acciones interrelacionadas dispuestas en un contexto temporal).

En el contexto del comercio minorista, el HAR enfrenta el desafío de discriminar entre interacciones comerciales legítimas, tales como examinar la etiqueta de un producto, comparar artículos o depositarlos en una canasta de compra, e interacciones anómalas o potencialmente delictivas, como el ocultamiento rápido de un objeto en prendas de vestir o bolsas personales.

La consecución de esta tarea requiere articular descriptores espaciales que capten la postura y la interacción mano-objeto con modelos temporales que analicen la progresión de dicha conducta en el tiempo \cite{chinthulla_deep_2025}.


\subsection{Modelado Conductual Pre-Delictivo}

\subsubsection{El Marco del Comportamiento Pre-Delictivo (\textit{Pre-Crime Behavior})}

La detección convencional de delitos mediante aprendizaje profundo ha tendido a enfocarse en la clasificación del acto ilícito en el instante preciso de su consumación.

Sin embargo, en el ámbito del hurto minorista, señalar el delito únicamente cuando el producto ya ha sido ocultado reduce el margen de acción del personal de seguridad, transformando la alerta automatizada en una confirmación tardía de la pérdida.

Para resolver esta falencia, Martínez-Mascorro et al. \cite{martinez-mascorro_suspicious_2021} formalizaron el marco del Comportamiento Pre-Delictivo (\textit{Pre-Crime Behavior} -- PCB), un enfoque conceptual que reinterpreta el análisis de video al modelar la secuencia de conductas preliminares que preceden a la ejecución del robo.

El marco PCB postula que la comisión de un hurto no es un evento instantáneo e aislado, sino el resultado de un proceso conductual estructurado en etapas temporales bien definidas. Dichas fases comprenden:

\begin{enumerate}
    \item Momento de Primera Aparición (\textit{First Appearance Moment} -- FAM): representa el ingreso y desplazamiento natural del individuo en el espacio comercial.

    \item Momento Comprensivo del Crimen (\textit{Comprehensive Crime Moment} -- CCM): constituye la etapa pre-delictiva propiamente dicha, caracterizada por acciones sutiles de preparación o exploración, tales como la verificación constante del entorno, giros de cabeza repetitivos para ubicar cámaras o personal, y la manipulación inusual de productos o prendas.

    \item Momento Estricto del Crimen (\textit{Strict Crime Moment} -- SCM): corresponde al acto físico del hurto o apoderamiento del objeto sin intención de pago.

    \item Momento Posterior al Hurto (\textit{Post-Theft Moment} -- PTM) y Retorno a la Normalidad (\textit{Back to Normality Moment} -- B2NM): fases en las cuales el sujeto busca disimular la acción y reincorporarse al flujo ordinario de clientes para abandonar el establecimiento.
\end{enumerate}

La ventaja teórica y práctica de adoptar el marco PCB radica en desplazar la frontera de detección hacia la fase CCM.

Al entrenar modelos de aprendizaje profundo capaces de reconocer los patrones de movimiento característicos del comportamiento sospechoso previo al robo, el sistema de videovigilancia genera alertas tempranas que permiten la intervención disuasiva antes de que la mercancía sea sustraída \cite{chinthulla_deep_2025}.


\subsection{Arquitecturas de Aprendizaje Profundo Espacio-Temporales}

\subsubsection{Extracción de Características Espaciales}

El procesamiento analítico de flujos de video requiere, en primera instancia, transformar la información pictórica no estructurada de cada fotograma en representaciones vectoriales densas que condensen la semántica de la escena.

Las Redes Neuronales Convolucionales (\textit{Convolutional Neural Networks} -- CNN) representan el pilar fundamental para esta tarea. A través de operaciones de convolución, agrupación (\textit{pooling}) y funciones de activación no lineales, las CNN extraen de forma jerárquica características visuales de bajo nivel (bordes, texturas y gradientes de color) hasta abstraer patrones de alto nivel, tales como siluetas humanas, regiones de interés o categorías de productos \cite{anikwe-onyiwa_recent_2025}.

En tareas que exigen la localización precisa de sujetos en tiempo real, los detectores de una sola etapa (\textit{single-stage detectors}) de la familia YOLO (\textit{You Only Look Once}) han reemplazado a las arquitecturas tradicionales de dos etapas debido a su alta eficiencia computacional.

Al dividir la imagen de entrada en una grilla regular y predecir simultáneamente los recuadros delimitadores (\textit{bounding boxes}) y las probabilidades de clase, algoritmos como YOLOv8 logran clasificar y localizar personas u objetos a altas tasas de fotogramas por segundo \cite{hameed_real_2025}.

Sin embargo, como señalan Pavithra et al. \cite{v_deep_2024} y Solórzano y Candela \cite{aguero_intelligent_2025}, los detectores espaciales puros operan sobre fotogramas de manera aislada, careciendo de la dimensión de memoria necesaria para distinguir si un movimiento identificado en una imagen fija constituye un gesto de compra legítimo o el inicio de una maniobra de ocultamiento.


\subsubsection{Modelado de Secuencias Temporales y Memoria Recurrente}

Dado que las conductas analizadas bajo el marco PCB se despliegan a lo largo del tiempo, la extracción de características espaciales debe complementarse con modelos capaces de procesar secuencias temporales.

Las Redes Neuronales Recurrentes (\textit{Recurrent Neural Networks} -- RNN) introducen conexiones con realimentación que permiten mantener un estado oculto interno, actuando como una memoria del contexto pasado.

No obstante, las RNN convencionales sufren del problema de desvanecimiento o explosión del gradiente al procesar secuencias prolongadas, lo que limita su capacidad para aprender dependencias a largo plazo.

Para superar esta restricción, Hochreiter y Schmidhuber introdujeron las redes de Memoria a Corto y Largo Plazo (\textit{Long Short-Term Memory} -- LSTM).

Una unidad LSTM incorpora una celda de memoria regulada por tres mecanismos de compuerta:

\begin{itemize}
    \item La compuerta de olvido (\textit{forget gate}), que decide qué información previa descartar.

    \item La compuerta de entrada (\textit{input gate}), que selecciona qué nuevos datos almacenar en el estado de la celda.

    \item La compuerta de salida (\textit{output gate}), que determina el valor emitido hacia el siguiente paso temporal.
\end{itemize}

La integración funcional entre la extracción espacial convolucional y el modelado temporal recurrente se consolida en la arquitectura de Redes Convolucionales Recurrentes de Largo Plazo (\textit{Long-Term Recurrent Convolutional Network} -- LRCN).

En una red LRCN, una serie de fotogramas continuos es procesada individualmente por un módulo CNN (\textit{TimeDistributed}), el cual genera una secuencia de vectores de características. Esta secuencia es posteriormente canalizada hacia capas LSTM que evalúan la evolución temporal de la conducta.

Como demostró Chinthulla \cite{chinthulla_deep_2025}, la arquitectura LRCN resulta especialmente apta para el análisis del hurto comercial, puesto que la componente convolucional captura la interacción visual entre el cliente y el entorno, mientras que la componente LSTM procesa la transición suave entre las etapas del comportamiento pre-delictivo.


\subsection{Edge Computing y Eficiencia Computacional}

\subsubsection{Procesamiento en Edge Computing}

El paradigma tradicional de analítica de video ha dependido de la transmisión de flujos de video continuos hacia servidores centralizados en la nube o estaciones de trabajo de alto rendimiento equipadas con GPUs dedicadas.

No obstante, en el contexto del comercio minorista de pequeña y mediana escala (minimarkets y bodegas), esta arquitectura centralizada resulta inviable debido a los costos recurrentes de ancho de banda, la latencia de red en la transmisión y la falta de infraestructura de cómputo avanzada \cite{aguero_intelligent_2025}.

El edge computing plantea descentralizar el procesamiento, ejecutando los algoritmos de inferencia directamente en dispositivos de hardware locales integrados al punto de captura de video.

Al procesar la información en tarjetas de cómputo embebido de bajo consumo, tales como las plataformas NVIDIA Jetson o Raspberry Pi, se elimina la dependencia de la nube, se reduce la latencia de respuesta a milisegundos y se resguardan las garantías de privacidad al evitar la transferencia externa de imágenes de videovigilancia \cite{hameed_real_2025, pranav_deep_2023}.


\subsubsection{Costo Computacional y Métricas de Inferencia}

El despliegue de modelos de aprendizaje profundo en dispositivos de borde impone restricciones estrictas respecto a la complejidad algorítmica y el consumo de recursos.

El costo computacional de un modelo se evalúa formalmente mediante indicadores de eficiencia tanto estructurales como de inferencia en tiempo de ejecución:

\begin{enumerate}
    \item Volumen de parámetros: determina la cantidad de pesos entrenables de la red y su huella de memoria RAM/VRAM en el dispositivo local.

    \item Latencia de inferencia: representa el tiempo requerido (medido en milisegundos) por el sistema para procesar una entrada y emitir un diagnóstico de clasificación.

    \item Tasa de procesamiento (\textit{Frames Per Second} -- FPS): indica el número de fotogramas por segundo que el modelo es capaz de evaluar continuamente. Para considerarse operativo en tiempo real en sistemas CCTV, la tasa de procesamiento debe situarse en el rango de 15 a 30 FPS.
\end{enumerate}

Como argumentan Ansari y Singh \cite{ansari_optimized_2023}, la optimización del costo computacional en modelos espacio-temporales no solo depende de la ligereza de la red base, sino del ajuste preciso de hiperparámetros como la longitud de la secuencia (\textit{Sequence Length}) y el tamaño de lote (\textit{batch size}), los cuales dictan la ventana de memoria necesaria para evaluar la conducta sin saturar la capacidad de cómputo local.


\subsection{Métricas de Evaluación de Rendimiento}

\subsubsection{Métricas de Clasificación y Detección Espacio-Temporal}

La validación empírica de un sistema de detección de hurtos exige métricas matemáticas rigurosas que contemplen tanto la corrección en la clasificación conductual como la precisión en la localización espacial.

En problemas de clasificación binaria o multiclase donde las muestras de delitos son sustancialmente menores que los eventos de compra normales, la exactitud simple (\textit{Accuracy}) resulta insuficiente por verse distorsionada por el desbalance de clases \cite{rashvand_exploring_2025}.

Por ello, la evaluación se fundamenta en las siguientes métricas derivadas de la matriz de confusión:

\begin{enumerate}
    \item Precisión (\textit{Precision}): mide la proporción de alertas de sospecha correctamente identificadas respecto al total de alertas emitidas por el modelo, siendo el indicador crítico para evaluar la minimización de falsos positivos:

    \begin{equation}
        \text{Precisión} = \frac{TP}{TP + FP}
    \end{equation}

    donde $TP$ representa los verdaderos positivos y $FP$ los falsos positivos.

    \item Exhaustividad o Sensibilidad (\textit{Recall}): calcula la proporción de eventos delictivos reales detectados por el sistema sobre el total de hurtos efectivamente acontecidos, midiendo la capacidad para mitigar los falsos negativos ($FN$):

    \begin{equation}
        \text{Recall} = \frac{TP}{TP + FN}
    \end{equation}

    \item Puntuación F1 (\textit{F1-Score}): constituye la media armónica entre la Precisión y el Recall, ofreciendo una métrica de desempeño global balanceada en escenarios desproporcionados:

    \begin{equation}
        \text{F1-Score} =
        2 \times
        \frac{\text{Precisión} \times \text{Recall}}
        {\text{Precisión} + \text{Recall}}
    \end{equation}

    \item Precisión Media Promedio (\textit{Mean Average Precision} -- mAP): en modelos que incorporan detección de objetos o delimitación de sujetos, la mAP evalúa el área bajo la curva de precisión-recall en función del solapamiento espacial (\textit{Intersection over Union} -- IoU) entre el recuadro predicho y la anotación de referencia.

    Se evalúa principalmente mediante mAP50 (con un umbral IoU de 0.50) y mAP50-95 (promediado en umbrales de IoU desde 0.50 hasta 0.95 en incrementos de 0.05), siendo esta última la prueba más estricta sobre la solidez del encuadre visual bajo situaciones de oclusión parcial \cite{aguero_intelligent_2025}.
\end{enumerate}
